\documentclass{jsarticle}
\usepackage{amsmath,amssymb}

%!!!! 重要 !!!! 重要 !!!! 重要 !!!! 重要 !!!! 重要 !!!! 重要 !!!! 重要 !!!!
%! 以下LaTeXスクリプト中、修正が必要な部分を（この文と同様に）%! で始まる
%! コメント文として提示します。修正を指示された部分をすべてコメント文中の
%! 指示に従って直してください。
%! - 修正は指示のある部分だけとし、スクリプト全体を新たに書き直したりしないこと.
%! - スクリプトの雛形に含まれている修正箇所の指示の文言を削除しないこと。
%! この指示に従わない場合、再提出とすることがあります。
%! (文章に誤植等があったとしても、指示された場所以外を修正する必要はありません。)

%% まずはこのスクリプトを一旦コンパイルしてPDFに変換してみるとよい．
%% このスクリプトは不完全ですが，コンパイル自体は成功するようになっています．

%% 【スクリプトのヒント】
%% - 数学モード中、\vec{式}で式の上に→をつけたベクトルを組版します．
%% - プランク定数ℏは数学モードでコマンド \hbar で組版します．

\begin{document}
\title{スピンとは何か --- 発見からディラック方程式まで\thanks{%
    藤川 和男 著「スピンとは何か --- 発見からディラック方程式まで」，
    数理科学 (サイエンス社) 2021年7月号，pp.~6--14 掲載
    より抜粋・改編．}}
\author{} %%% 著者を表示しない
\date{}   %%% 日付を表示しない
\interfootnotelinepenalty=10000 % 脚注が最初のページにうまく表示されるようにおまじない

\maketitle

\section{パウリの排他律}
\label{sec:pauliexclusion}

%!!!!!!!!!!!!!!!!!!!!!!!!!!!!!!!!!!!!!!!!!!!!!!!!!!!!!!!!!!!!!!!!!!!!!!!!!!!!!
%!q1. この節の2段落目の組版がおかしい。(1行目の行頭が下がっていない。)
%! LaTeX資料3.2節14ページの「強制改行」に関する注意書きをよく読んで正しく
%! 段落が組版されるように修正すること。

量子力学の発見以前の前期量子論の段階で，量
子物理学の多くの基本的な法則等が発見されたが，
本稿では電子に関する量子数とその特性に関する
発見の歴史を述べたい．\\
水素原子の基本的な性質の説明で成功を収めたボーアを中心として
周期律表の説明が試みられた．よく知られているように
ボーアは水素原子の模型において，3つの量子数,
主量子数$n$, 角運動量$l$, その$z$成分$m$
\[
  n, l, m (=l_z)
\]
を仮定し，$m=l，l-1，\ldots，-l$とした．
これを一般化して，ある種の周期性は理解できたが，
周期律表の $2,8,18,\ldots$ という数は説明できなかった.
また磁場中の原子からのスペクトル線の分離であるゼーマン効果の説明も，軌道運動に関係した
正常ゼーマン効果は荷電粒子の円運動で説明できたが，古典的な対応のない異常ゼーマン効果は説明できなかった.

このような中で，当時25歳のイギリスのストーナー(E.\ Stoner)は，次のような
追加の自由度を含む4個の量子数を用いると多くの事実が説明できることに気づいた：
\[
  n, l, m (=j_z)
\]
ここで
\[
  j = l\pm \frac{1}{2}
\]
である．この考えはゾンマーフェルトに評価され，パウリにも採用された．より精密な
分析がパウリにより与えられ，特に$2,8,18,\ldots$ という数の説明，原子の安定性
(すべての電子が基底状態に落ち込まないこと) を説明する方策として，これらの
4個の量子数で指定される状態には，2個以上の電子は入ることが許されないという
「排他律」の考えを提案した．異常ゼーマン効果も，通常のゼーマン効果で知られた
磁気モーメントと磁気量子数の関係 
%!!!!!!!!!!!!!!!!!!!!!!!!!!!!!!!!!!!!!!!!!!!!!!!!!!!!!!!!!!!!!!!!!!!!!!!!!!!!!
%!q2. 以下の括弧に囲まれた文中の数式(4箇所)を埋めよ。
%!    いずれの数式も$...$の形の数学モードの数式として組版すること。
(以降では，混同する可能性がある場合は，
電子質量は $m_e$ と記す．また電子の電荷$q$は$e>0$を用いて$q=-e$とする．) 
\[
    \frac{e\hbar}{2m_{e} c} m 
\]
($m=-1,0,1$) に対応して，新しい古典的な対応のない自由度に関しては
\[
    2\times \frac{e\hbar}{2m_{e} c} m \times \left(\pm\frac{1}{2}\right)
\]
と取れば，説明できることが示された．この最初
の因子 2 (これは後にトーマス因子と呼ばれること
になった) がその後の研究において重要となった．


\section{スピン角運動量の提案}
\label{sec:angularmomentum}

こで現れた$\pm\frac{1}{2}$という新しい量子数で表され
る古典的な対応のない二価性の本質は何なのかが
問題になるが，1925年当時はパウリ等はより深い
量子現象の理解が得られたときに解明できるも
のと考えた．これに反して，この二価性を電子が持つコマのような
回転に起因する固有な角運動量と考える提案が，アメリカから来た当時20歳の
クローニッヒ (R.\ Kronig) により考えられたが，意見を
聞いたパウリ，ハイゼンベルグ，ボーア等が否定的な
意見を述べ，発表をためらった．ちなみに，パウリもハイセンベルグも当時は25歳で
あった．そうこうするうちに，ライデンのエーレンフェストのところにいた
若い研究者のガウシュミット (G.\ Goudsmit) および
ウーレンベック (G.\ Uhlenbeck) の 2人が，この二価性はコマのように
回転する電子の自由度に起因するという
アイデアの論文を書いた．このときも， 2人は，パウリとか
ローレンツ等に話したところ非常に不評だったので，先生のエーレンフェストの
ところへ出かけてあの論文を取り下げたいと
申し出たところ，エーレンフェストは ``あの論文はすでに
雑誌に投稿した．君たちは若いから馬鹿な論文を書いても失うものがない'' 
と言ったといわれている．

しかし，時が経つにつれてスピンは本当に回転
しているか否かは別にして，質量とか電荷と同じように電子が持つ固有な
角運動量であるという考えが多くの人たちの支持を
得るようになった．またこのころ (1925年から26年にかけて)，ハイゼンベルグとか
シュレーディンガーにより量子力学が提案された．このような発展を受けて，パウリ
はスピンの実体論的な理解を追及ぜずに，その数学的な記述を与えた．

%!!!!!!!!!!!!!!!!!!!!!!!!!!!!!!!!!!!!!!!!!!!!!!!!!!!!!!!!!!!!!!!!!!!!!!!!!!!!!
%!q3.  次の行中、太字になるべきところがそうなっていない。修正せよ。
今日パウリの\textbf{スピン行列}と呼ばれている次の行列を定義した：
%!!!!!!!!!!!!!!!!!!!!!!!!!!!!!!!!!!!!!!!!!!!!!!!!!!!!!!!!!!!!!!!!!!!!!!!!!!!!!
%!q4.  以下の式中の行列表示を完成させよ。行列表示にはpmatrix環境を用いること。
%!     (LaTeX資料4.2.3節参照) 
\begin{equation}
    \sigma_1 = \begin{pmatrix}
        0 & 1 \\
        1 & 0
    \end{pmatrix},
    \sigma_2 = \begin{pmatrix}
        0 & -i \\
        i & 0
    \end{pmatrix},
    \sigma_3 = \begin{pmatrix}
        1 & 0 \\
        0 & -1
    \end{pmatrix}.
    \label{eq:pauliMatrices}  %! この行は削除しないこと （数式に対する相互参照用のラベルを定義しています）
\end{equation}
これらのスピン行列は交換関係および反交換関係
\[
[\sigma_1,\sigma_2] =2i\sigma_3,
[\sigma_2,\sigma_3] =2i\sigma_1,
[\sigma_3,\sigma_1] =2i\sigma_2,
\{ \sigma_k,\sigma_l \} = 2\delta_{kl}
\]
を満たす．スピン角運動量を
\[
s_x = \frac{\hbar}{2} \sigma_1,
s_y = \frac{\hbar}{2} \sigma_2,
s_z = \frac{\hbar}{2} \sigma_3,
\]
で定義すると，角運動量の交換関係
\[
[s_x,s_y] = i\hbar\sigma_z,
[s_y,s_z] = i\hbar\sigma_x,
[s_z,s_x] = i\hbar\sigma_y
\]
を満たし，角運動量の大きさは$s=\frac{1}{2}$に対応する
\[
\vec{s}^2 = s_x^2 + s_y^2 + s_z^2 = \hbar^2\frac{1}{2}(\frac{1}{2}+1) 
    \begin{pmatrix}
    1 & 0 \\
    0 & 1    
    \end{pmatrix}
\]
である．最後の行列は単位行列を表す．スピン自由度を表す波動関数としては$s_z$の固有関数として
\[
s_z\alpha = \frac{\hbar}{2}\alpha,\vec{s}^2\alpha = \hbar^2\frac{3}{4}\alpha,
s_z\beta = -\frac{\hbar}{2}\beta,\vec{s}^2\beta = \hbar^2\frac{3}{4}\beta
\]
を満たす2成分の直交ベクトルが選ばれる．全角運動量は
\[
\vec{j} = \vec{l} + \vec{s}
\]
で与えられる.

スピンの自由度を含まないシュレーディンガー方程式を
%!!!!!!!!!!!!!!!!!!!!!!!!!!!!!!!!!!!!!!!!!!!!!!!!!!!!!!!!!!!!!!!!!!!!!!!!!!!!!
%!q5. 以下の式を完成させよ。ただし、{ } や [ ] の括弧を大きく組版(LaTeX資料 4.2.5節参照)すること。
\[
i\hbar\frac{\partial}{\partial t}\psi(t,\vec{x}) =
\left\{
\frac{1}{2m_e} \left[ \frac{\hbar}{i}\vec{\nabla} + \frac{e}{c}\vec{A}(t,\vec{x}) \right]^2 - e A_0(t,\vec{x})
\right\}\psi(t,\vec{x})
\]
の形に書くときは，スピンを含めた解は直積の形で書かれた
$\alpha\psi(t,\vec{x})$
と
$\beta\psi(t,\vec{x})$
およびそれらの線形結合で与えられる．


\section{パウリ方程式}
\label{sec:paulieq}

スピンの自由度と通常の時空の自由度を直積ではなく，もう少し
本質的な形で表す式として，パウリは次のシュレーディンガー方程式の
一般化(\textbf{パウリ方程式})を提案した：
%!!!!!!!!!!!!!!!!!!!!!!!!!!!!!!!!!!!!!!!!!!!!!!!!!!!!!!!!!!!!!!!!!!!!!!!!!!!!!
%!q6. 以下の式をalign環境(LaTeX資料 4.2.6節参照)を用いて数式を組版せよ。（環境を変更しないこと）
% その際、等号=の位置を揃えること。また、数式の{ } や [ ] の括弧を大きく組版(LaTeX資料 4.2.5節参照)すること。
\begin{align}
    i\hbar\frac{\partial}{\partial t}\psi(t,\vec{x}) 
    = &
    \left\{
        \frac{1}{2m_e} \left[ \left(\frac{\hbar}{i}\vec{\nabla} + \frac{e}{c}\vec{A}(t,\vec{x})\right)\cdot\vec{\sigma} \right]^2 - e A_0(t,\vec{x})
    \right\}\times\psi(t,\vec{x})  
    \notag\\ %%! この行は削除してはいけない （\notag でこの行については数式番号を出力しないように指定しています）
    = & \label{eq:paulieq} %! この行は削除しないこと （数式に対する相互参照用のラベルを定義しています）
    \left\{
        \frac{1}{2m_e} \left[ \frac{\hbar}{i}\vec{\nabla} + \frac{e}{c}\vec{A}(t,\vec{x}) \right]^2 - e A_0(t,\vec{x})
    \right\} -2\times \frac{\hbar e}{2m_e c}  \vec{s}\cdot\vec{B}\psi(t,\vec{x}) 
\end{align}
波動関数は2成分の関数
%!!!!!!!!!!!!!!!!!!!!!!!!!!!!!!!!!!!!!!!!!!!!!!!!!!!!!!!!!!!!!!!!!!!!!!!!!!!!!
%!q7. 以下の式を完成させよ。ベクトルの成分表示はpmatrix環境を用いて2×1行列として組版せよ。
\[
\psi(x) = \begin{pmatrix}
    \psi_1(x) \\
    \psi_2(x)  
\end{pmatrix}
\]
で与えられる．式\eqref{eq:paulieq}の最後の項が磁場と電子の
スピンに起因する磁気モーメントの相互作用を表しており，トーマス因子を
正しく出している．このようにして，非相対論的な近似での
電子とそのスピン自由度を正しく記述する定式化が
与えられた．1927年のことである．

電子のスピンのアイデアの歴史を
振り返ると，今日の視点で考えて興味深いのは，パウリをはじめ当時の
中心的な人たちは，ボーアの対応原理では理解できない (古典的な対応物のない) 
角運動量 1/2 (あるいは二価性) を理解していたが，それを電子の
固有な性質 (回転するコマのようなもの) という実体論的な解釈には
反対したことである．もっとも，後の量子力学とか場の量子論という観点から見ても
スピンの実体論的な理解は現在でも進歩が見られていないことから，当時の
人たちが慎重であったことも本質を掴んでいたともいえる．特に
パウリの抽象的な2行2列の行列を用いた表示は本質をついていた．



\section{ディラック方程式}
\label{sec:diraceq}

パウリによるスピン行列およびパウリ方程式は，多くの
応用をもたらした．しかし，相対論的な不変性は満たされていない．ディラックは
この欠点を正し，追加の相対論的な重要な効果を明らかにしたうえ，
さらに反粒子の存在という驚くべき発見に導いた．クライン・ゴルドン
方程式は，アインシュタインの関係式$p_{\mu}^2 - m^2 =0$を波動関数に作用させて
$[(i\hbar\partial_{\mu})^2 - m^2]\phi(x)  =0$の形の相対論的な方程式を与えるが，この式は時間に
関して2回の微分を含んでおり，確率密度と同定されるものが正定値ではないので
シュレーデインガー方程式の相対論的不変な一般化と見なせない．

%!!!!!!!!!!!!!!!!!!!!!!!!!!!!!!!!!!!!!!!!!!!!!!!!!!!!!!!!!!!!!!!!!!!!!!!!!!!!!
%!q8. 以下の式を完成させよ。
ディラックはこの困難を解決するために，時間に関して1階の相対論的不変な方程式を
探求した．模式的には，これはアインシュタインの関係式$g^{\mu\nu}p_{\mu}p_{\nu} - (mc)^2=0$の
ルートを考えることに対応していた．すなわち，
\[
\sqrt{g^{\mu\nu}p_{\mu}p_{\nu} - (mc)^2} \sim \gamma^\mu p_{\mu} - (mc)
\]
とおいて見ることに対応する．この逆の操作を考えると
%!!!!!!!!!!!!!!!!!!!!!!!!!!!!!!!!!!!!!!!!!!!!!!!!!!!!!!!!!!!!!!!!!!!!!!!!!!!!!
%!q9. 以下の式をalign*環境(LaTeX資料 4.2.6節参照)を用いて数式を組版せよ。（環境を変更しないこと）
% その際、等号=の位置を揃えること。
\begin{align*}
     & [\gamma^{\mu} p_{\mu} +(mc)] [\gamma^{\nu} p_{\nu} - (mc)] \\
   = & \gamma^{\mu} \gamma^{\nu} \frac{1}{2}(p_{\mu}p_{\nu} + p_{\nu} p_{\mu}) - (mc)^2 \\
   = & \frac{1}{2} (\gamma^{\mu} \gamma^{\nu} +  \gamma^{\nu} \gamma^{\mu}) p_{\mu}p_{\nu} - (mc)^2 \\
   = & g^{\mu\nu} p_{\mu}p_{\nu} - (mc)^2
\end{align*}
を求めることになる.
$[\gamma^{\nu}\hat{p}_{\nu} - (mc)]\psi(x) = 0$が成立すれば
\[
[\gamma^{\nu}\hat{p}_{\nu} + (mc)][\gamma^{\nu}\hat{p}_{\nu} - (mc)]\psi(x) 
= [g^{\mu\nu} \hat{p}_{\mu}\hat{p}_{\nu} - (mc)^2]\psi(x) = 0
\]
が言えるからである．この考察から
%!!!!!!!!!!!!!!!!!!!!!!!!!!!!!!!!!!!!!!!!!!!!!!!!!!!!!!!!!!!!!!!!!!!!!!!!!!!!!
%!q10. 以下の式は添字の組版が少しおかしい。修正せよ。
\[
\gamma^\mu \gamma^\nu +  \gamma^\nu \gamma^\mu = 2 g^{\mu\nu}
\]
が要求される．この関係式は普通の数では満たされない．

ディラックは4行4列の行列$\gamma^{\mu}$がこの関係式
を満たすことを示した．具体的には2行2列の
%!!!!!!!!!!!!!!!!!!!!!!!!!!!!!!!!!!!!!!!!!!!!!!!!!!!!!!!!!!!!!!!!!!!!!!!!!!!!!
%!q11.  以下の行で相互参照用のラベルが正しく設定されていないため、次の段落冒頭の
%!   式番号が正しく組版されない。LaTeX資料3.1.3節「相互参照」を読み、
%!   適切なラベル名を設定することで式番号が正しく組版されるように修正すること。
%!   (必ず相互参照の機能を使うこと。式番号をそのまま(3)のように書かないこと。)
パウリ行列\eqref{eq:pauliMatrices}を用いて次の4行4列の表示で与えられる：
\[
\gamma^0 = \begin{pmatrix}
    1 & 0 \\ 0 & -1
\end{pmatrix},
\gamma^k = \begin{pmatrix}
    0 & \sigma_k \\ -\sigma_k & 0
\end{pmatrix},
k=1,2,3.
\]
この表示では$\gamma^0$の$1$および$0$はそれぞれ2行2列の単位行列および
ゼロ行列を表す．ディラック理論で基本的となるカイラルな性質を記述する$\gamma_5$行列は
\[
\gamma_5 = i \gamma^0 \gamma^1 \gamma^2 \gamma^3,
\gamma_5 \gamma^{\mu} + \gamma^{\mu} \gamma_5 = 0
\]
で定義され，現在の$\gamma$行列の表示では具体的に4行4列の
\[
\gamma_5 = \begin{pmatrix}
    0 & 1 \\ 1 & 0
\end{pmatrix},
(\gamma_5)^2 = \begin{pmatrix}
    1 & 0 \\ 0 & 1
\end{pmatrix} 
\]
で与えられる．
また，4行4列のディラック行列
$\gamma^{\mu}$,
$\mu = 0,1,2,3$は
\[
(\gamma^0)^{\dagger} = \gamma^0,
(\gamma^k)^{\dagger} = -\gamma^k,
k=1,2,3
\]
を満たす．例えば，$(\gamma^0)^{\dagger}$は$\gamma^0$のエルミート共役
を表す．

こうして，ディラックは4成分を含む$\psi(x)$に対する\textbf{ディラック方程式}
\begin{equation} \label{eq:diraceq}
    [i\hbar \gamma^{\mu} \partial_{\mu} - mc]\psi(x) = 0
\end{equation}
を発見した．1928年のことである．

\end{document}